% Chapter 7: Algebraic Cipher Types
% DEFINES: ch:algebraic-types; sec:type-algebra, sec:product-types,
%   sec:sum-impossibility, sec:granularity, sec:ch7-notes;
%   prop:product-tradeoff, thm:sum-impossibility, ex:product-pair,
%   rem:granularity, tab:constructors
% REFERENCES (backward): sec:cipher-type, rem:values-are-maps, def:uniformity,
%   def:cipher-map
% REFERENCES (forward, part labels): part:confidentiality, part:constructions

\chapter{Algebraic Cipher Types}
\label{ch:algebraic-types}

\Cref{sec:cipher-type} named the cipher type $\cipher{X}$ and promised an
algebra. This chapter delivers it, and the delivery is uneven in a way that
turns out to be the point. The type constructors of ordinary programming,
products, sums, functions, do not all survive the passage through the
trapdoor equally. Products pass, at a cost one can choose; functions are the
cipher maps themselves; but sums hit a genuine impossibility, a wall no
construction climbs. The chapter is organized around that asymmetry, because
where a structure fails is more informative than where it succeeds.

\section{Cipher Types as a Type Algebra}
\label{sec:type-algebra}

A cipher type $\cipher{X}$ is the type of cipher values over $X$, and by
\Cref{rem:values-are-maps} its inhabitants are themselves cipher maps from a
trivial domain. The question of this chapter is which constructions on
types lift to cipher types: given $\cipher{A}$ and $\cipher{B}$, is there a
sensible $\cipher{A \times B}$, a $\cipher{A + B}$, a $\cipher{A \to B}$, and
what does each cost?

The two trivial types set the floor. The void type $\mathbf{0}$ has no
values, so a cipher map out of it is vacuously correct and carries nothing;
to the untrusted machine it is a random function never validly queried. The
unit type $\mathbf{1}$ has a single value, so a cipher value over it is a
cipher constant, and totality keeps even that single value hidden among
noise. Neither is interesting on its own, but they anchor the algebra.
\Cref{tab:constructors} is the chapter's map: it states, for each
constructor, how it lifts and what the lift costs, and the sections below
fill it in.

\begin{table}[t]
\centering
\caption{The type constructors and how they lift to cipher types.}
\label{tab:constructors}
\begin{tabular}{@{}lll@{}}
\toprule
Constructor & Lifts to a cipher type? & Cost or obstruction \\
\midrule
Void $\mathbf{0}$ & yes, vacuously & none \\
Unit $\mathbf{1}$ & yes, a cipher constant & none \\
Product $A \times B$ & yes & projections need a cipher map, or the joint leaks \\
Sum $A + B$ & not fully & tag-hiding and untrusted dispatch are exclusive \\
Function $A \to B$ & yes, it is a cipher map & the framework itself \\
\bottomrule
\end{tabular}
\end{table}

The function type is already in hand: a cipher map for $f : A \to B$
\emph{is} the cipher exponential $\cipher{A \to B}$, the trapdoor object the
whole book studies. That leaves the two data constructors, product and sum,
and they are where the algebra becomes informative.

\section{Product Types Pass}
\label{sec:product-types}

A product can be encoded two ways, and the choice is a confidentiality knob
rather than a wall.

Encode the pair $(a, b)$ \emph{jointly}, as one cipher value over
$A \times B$, and the untrusted machine sees a single opaque token: it
cannot read $a$ or $b$, and it cannot even detect a correlation between
them, since the correlation is buried inside one encoding. The price is
that the projections are gone. Recovering $a$ from a joint encoding of
$(a, b)$ is itself a function on tokens, hence a cipher map the trusted
machine must build; the product hides its components by making them
inaccessible without a constructed projector. Encode the pair
\emph{component-wise} instead, as a token for $a$ beside a token for $b$,
and the projections are free, just take one side, but the joint distribution
is no longer hidden: each component is a deterministic function of its
value, so an adversary who watches enough pairs reconstructs how $a$ and $b$
covary, even when each component alone is near-uniform.

\begin{proposition}[Product confidentiality trade-off]
\label{prop:product-tradeoff}
For finite types $A$ and $B$:
\begin{enumerate}
\item Under joint encoding $\cipher{A \times B}$, the untrusted machine
  cannot distinguish $(a_1, b_1)$ from $(a_2, b_2)$ beyond the
  representation-uniformity advantage $\delta$, but each projection requires
  a cipher map the trusted machine constructs.
\item Under component-wise encoding $\cipher{A} \times \cipher{B}$,
  projections are free, but an adversary observing $N$ pairs can estimate
  the joint distribution of $(A, B)$ to total-variation accuracy $\xi$ once
  $N = O(\lvert A\rvert\,\lvert B\rvert / \xi^2)$.
\end{enumerate}
\end{proposition}

The proof of part~(1) is the observation that a projection is a token-to-token
map, hence a cipher map; part~(2) is a counting argument on the empirical
joint distribution of cipher pairs, given in \cite{towell2026algebraic}. The
point for the algebra is that the product \emph{passes}: a usable
$\cipher{A \times B}$ exists either way, and the two encodings are the two
ends of a dial between hidden correlations and free projections. Nothing is
impossible here; one only chooses.

\begin{example}[A pair, two ways]
\label{ex:product-pair}
Let $A = \{a_1, a_2\}$ and $B = \{b_1, b_2\}$, and suppose the latent pairs
are perfectly correlated, $a_1$ always with $b_1$ and $a_2$ with $b_2$.
Encoded jointly, the four-token cipher type $\cipher{A \times B}$ shows the
untrusted machine only that some pair occurred; the correlation $a_i
\leftrightarrow b_i$ is invisible, sealed inside each token. Encoded
component-wise, the same correlation is laid bare: token-for-$a$ and
token-for-$b$ co-occur in exactly two combinations out of the four possible,
and an adversary counting $N = O(\lvert A\rvert\,\lvert B\rvert/\xi^2)$ pairs
reads off the perfect correlation. Same data, same per-component uniformity,
opposite joint confidentiality, the granularity dial of
\Cref{sec:granularity}.
\end{example}

\section{The Sum-Type Impossibility}
\label{sec:sum-impossibility}

The sum type is where the algebra breaks, and the break is structural, not a
missing construction. A value of type $A + B$ carries a tag, left or right,
and ordinary pattern matching reads the tag to choose a handler. Two things
one would want from a cipher sum turn out to be mutually exclusive.

\begin{theorem}[Sum-type impossibility]
\label{thm:sum-impossibility}
For non-empty finite types $A$ and $B$, these two properties cannot both
hold of a cipher type for $A + B$:
\begin{enumerate}
\item \textbf{Tag hiding}: from a cipher value alone, the untrusted machine
  cannot tell whether it encodes an $A$-value or a $B$-value, beyond the
  uniformity advantage $\delta$.
\item \textbf{Untrusted dispatch}: holding the two handlers separately, the
  untrusted machine can apply exactly one of them to a cipher value,
  choosing which as a function of that value, with no per-query help from the
  trusted machine.
\end{enumerate}
\end{theorem}

The obstruction is clean once seen. Suppose both held. Untrusted dispatch
means a selector $s$ that reads a token and returns left or right, correctly
routing $A$-encodings to the $A$-handler and $B$-encodings to the
$B$-handler. But such an $s$ is exactly a test that separates $A$-encodings
from $B$-encodings, which is to say a tag distinguisher; if it routes
correctly it reads the tag, and tag hiding fails. The two requirements are
the same capability viewed from two sides: dispatching on the tag and
reading the tag are one act. Joint encoding takes the first horn, it hides
the tag and forgoes untrusted dispatch (the case analysis can still be
performed, by a fused eliminator that evaluates to a tagged-hidden result,
but the untrusted machine never forms the selector); component-wise encoding
takes the second, it exposes the tag in the clear and dispatches freely. No
encoding sits between.\footnote{The full proof, including the fused-eliminator
escape that produces the case result without forming the selector, is in
\cite{towell2026algebraic}.}

This is the chapter's sharpest contrast. The product (\Cref{sec:product-types})
offered a dial: both encodings give a working type, trading projections
against joint hiding. The sum offers a wall: the two things one wants are
provably incompatible, and a designer must give one up. The reason traces
to totality. Because every token decodes to something, a hidden tag leaves
the untrusted machine no handle to branch on, and any handle it could branch
on is a leak of the tag.

\section{Encoding Granularity and Entanglement}
\label{sec:granularity}

The product dial and the sum wall are two readings of one underlying
parameter: how much is encoded together. This is the place to pay off the
limitation \Cref{def:uniformity} flagged, that $\delta$ bounds only the
marginal distribution of a token and says nothing about joint distributions
across tokens.

Call the number of components encoded as a single unit the
\emph{entanglement} $p$. At $p = 1$, each value is encoded on its own; every
token can be made marginally near-uniform ($\delta \approx 0$), but
correlations between separately encoded values pass straight through, since
each token is a deterministic function of its value and an adversary can
build the joint empirical distribution. At the other extreme, $p = k$
encodes a $k$-tuple as one token, and the joint distribution of all $k$
components is hidden inside that single uniform-looking value, at a cost: the
token space must cover $\lvert Y\rvert^{k}$ possibilities, so space grows
exponentially in $p$. Between the extremes, $p$ is a dial from
marginal-only confidentiality to joint confidentiality, paid for in space.

\begin{remark}[The granularity principle]
\label{rem:granularity}
Coarser encoding (larger $p$) hides more correlations but blocks more
operations and costs more space; finer encoding (smaller $p$) frees
operations and saves space but leaks correlations. Property~2's $\delta$
lives entirely at the marginal scale, $p = 1$ per token; it is silent about
what \Cref{ex:product-pair} exposed. The joint distribution of two
separately encoded values is recoverable at total-variation accuracy $\xi$
from $O(\lvert Y_1\rvert\,\lvert Y_2\rvert / \xi^2)$ observations, a rate no
choice of $\delta$ changes. That rate, and its matching lower bound, is the
compositional confidentiality scale, governed by quantities $\delta$ does
not touch; \Cref{part:confidentiality} develops it, and \Cref{ch:orbit-closure}
gives its active-adversary face. The ideal of one giant joint encoding ($p$
maximal) is exactly the impractical extreme that composition exists to
avoid: one composes smaller cipher maps and accepts that intermediate
correlations may leak.
\end{remark}

\section{Notes and Provenance}
\label{sec:ch7-notes}

The type-constructor analysis, the product trade-off, and the sum-type
impossibility are the algebraic-cipher-types paper's, restated here with the
proofs sketched and pointed to the paper, per this part's proof-depth
decision. The granularity principle and the entanglement parameter $p$ are
the spine's (its encoding-granularity section), and the
$O(\lvert Y_1\rvert\lvert Y_2\rvert/\xi^2)$ joint-recovery rate is the bridge
to the compositional confidentiality of \Cref{part:confidentiality}, where it
acquires a matching lower bound. The sum-type impossibility is the structural
heart: it is the algebraic shadow of totality, the same property that bought
privacy in \Cref{ch:four-properties} now forbidding a hidden tag from being
branched on. The summary table (\Cref{tab:constructors}) is a candidate to
promote into the spine as the shared statement of which constructors lift.
